% Volume V: Dynamics, Signals, Graphs, and Control
% Continuity note for Chapter 29.
% Previous dependencies: Chapters 3--6 provide vectors, matrices, eigenvectors, quadratic forms, and PCA; Chapter 28 provides signal intuition and Fourier decompositions.
% New durable concepts: graphs as vertex-edge objects; adjacency matrices; degree; weighted and directed graphs; paths, neighborhoods, and connected components; graph signals; combinatorial and normalized Laplacians; Laplacian quadratic forms; random walks; diffusion on graphs; Laplacian eigenvectors; Fiedler vector; spectral clustering; graph Fourier transforms; network metrics; connectomes; functional networks; message passing and graph neural networks.
% Forward hooks: Chapter 30 uses graph-like state transition diagrams for dynamic programming; Chapter 36 uses transition graphs in reinforcement learning; Chapter 40 uses GNNs; Chapter 48 uses population geometry on networked neural data; Chapter 51 uses network summaries in neural data analysis.
% Notation commitments: Graphs are G=(V,E); |V|=n; adjacency matrix is A; degree matrix is D; Laplacian is L=D-A; graph signals are x\in\mathbb R^n; row random-walk matrix is P=D^{-1}A when degrees are positive.
\section{Chapter 29: Graphs, Networks, and Spectral Graph Theory}
\label{ch:graphs-networks-spectral-graph-theory}

Time series organize data along a line. Many neuroscience and AI problems organize data over networks instead: brain regions connected by axonal projections, neurons connected by synapses, genes linked by regulation, web pages linked by hyperlinks, papers linked by citations, or tokens connected by attention. The mathematical object is a graph.

\paragraph{Chapter problem.}
How do graphs represent relational structure, how do matrices encode that structure, and how do Laplacian eigenvectors reveal smoothness, diffusion, clusters, and modes on networks?

\paragraph{Section map.}
Section 29.1 defines graph basics, adjacency matrices, degrees, paths, components, random walks, and graph signals. Section 29.2 develops graph Laplacians, quadratic forms, smoothness, nullspaces, and diffusion. Section 29.3 uses Laplacian eigenvectors for graph Fourier analysis, spectral clustering, and embeddings. Section 29.4 connects network metrics to connectomes, neural data, and graph neural networks.

\subsection{29.1 Graph basics}

\paragraph{Guiding question.}
What are the basic objects needed to represent pairwise relationships among units?

\paragraph{Beginner-facing intuition.}
A graph is a set of objects together with connections among them. The objects are vertices or nodes. The connections are edges. The same graph idea can represent neurons and synapses, brain areas and tracts, users and friendships, molecules and bonds, or variables and dependencies.

Graphs are useful because the pattern of connections changes what computations or dynamics are possible. Two nodes with the same feature value can play different roles if one is a hub and the other is peripheral.

\paragraph{Formal definitions and notation.}
An undirected graph is
\[
G=(V,E),
\]
where \(V=\{1,\ldots,n\}\) is a finite vertex set and \(E\subseteq\{\{i,j\}:i\neq j\}\) is an edge set. A weighted undirected graph has nonnegative weights \(A_{ij}=A_{ji}\geq0\). The adjacency matrix is
\[
A\in\mathbb R^{n\times n},\qquad
A_{ij}=
\begin{cases}
w_{ij}, & \text{if connected},\\
0, & \text{otherwise}.
\end{cases}
\]
For a simple unweighted graph, \(A_{ij}\in\{0,1\}\) and \(A_{ii}=0\). The degree of node \(i\) is
\[
d_i=\sum_{j=1}^n A_{ij}.
\]
The neighborhood of \(i\) is
\[
N(i)=\{j:A_{ij}>0\}.
\]
A path is a sequence of vertices in which consecutive vertices are connected. A connected component is a maximal set of vertices joined by paths.

A graph signal is a vector
\[
x\in\mathbb R^n,
\]
where \(x_i\) is a value attached to node \(i\), such as activity of brain region \(i\), a label score for paper \(i\), or a molecular feature at atom \(i\).

\paragraph{Core mechanism: adjacency matrices count walks.}
For an unweighted graph, the matrix power \(A^k\) encodes walks:
\[
(A^k)_{ij}=\text{number of length-}k\text{ walks from }i\text{ to }j.
\]
\emph{Derivation.} For \(k=1\), \(A_{ij}\) is \(1\) exactly when there is a one-step walk. If \((A^k)_{i\ell}\) counts length-\(k\) walks from \(i\) to \(\ell\), then
\[
(A^{k+1})_{ij}=\sum_{\ell=1}^n (A^k)_{i\ell}A_{\ell j}.
\]
Each term counts a length-\(k\) walk from \(i\) to \(\ell\) followed by one edge from \(\ell\) to \(j\). Summing over \(\ell\) counts all possible last intermediate vertices.

\paragraph{Concrete example.}
Consider the path graph \(1-2-3\). Its adjacency matrix is
\[
A=
\begin{pmatrix}
0&1&0\\
1&0&1\\
0&1&0
\end{pmatrix}.
\]
The degrees are \(d=(1,2,1)\). The square is
\[
A^2=
\begin{pmatrix}
1&0&1\\
0&2&0\\
1&0&1
\end{pmatrix}.
\]
The entry \((A^2)_{13}=1\) counts the walk \(1\to2\to3\). The entry \((A^2)_{22}=2\) counts \(2\to1\to2\) and \(2\to3\to2\).

\paragraph{Computational neuroscience example.}
In a structural connectome, vertices may be brain regions and weighted edges may estimate fiber density between regions. A graph signal \(x_i\) could be average activity in region \(i\) during a task. The graph structure says which regions are anatomically close in the model; the signal says what is happening on those regions.

\paragraph{AI or machine-learning example.}
In a citation graph, nodes are papers and edges connect citations. A node feature may be a text embedding. A classifier can use both the paper's own features and its neighborhood. Graph neural networks formalize this by repeatedly aggregating features from neighboring nodes.

\paragraph{Common misconceptions and failure modes.}
A graph layout drawing is not the graph; many drawings can represent the same adjacency structure. Degree is local and does not by itself measure global importance. A directed graph need not have a symmetric adjacency matrix. A weighted edge must be interpreted in context: a high weight could mean strong anatomical projection, high correlation, physical proximity, or learned attention depending on the model.

\paragraph{Exercises.}
\begin{enumerate}[label=\textbf{29.1.\arabic*.}, leftmargin=*]
\item \textbf{Mechanical.} Write the adjacency matrix and degree vector for the four-node chain \(1-2-3-4\).
\item \textbf{Proof or derivation.} Prove by induction that \((A^k)_{ij}\) counts length-\(k\) walks in an unweighted graph.
\item \textbf{Numerical/computational.} For the triangle graph on three nodes, compute \(A^2\) and interpret the diagonal entries.
\item \textbf{Interpretation.} In a brain-region graph, what is the difference between a graph edge and a graph signal value?
\item \textbf{Misconception/debugging.} A student says, ``Node 5 is drawn in the center, so it must be the most connected.'' Explain why the drawing alone is not enough.
\end{enumerate}

\subsection{29.2 Graph Laplacians}

\paragraph{Guiding question.}
How can a matrix measure whether a graph signal changes smoothly across connected nodes?

\paragraph{Beginner-facing intuition.}
On a line, a derivative measures how rapidly a function changes with position. On a graph, there is no continuous coordinate, but there are edges. The graph Laplacian measures differences across edges. A graph signal is smooth if connected nodes have similar values.

This makes the Laplacian a central bridge between geometry, diffusion, clustering, and regularization.

\paragraph{Formal definitions and notation.}
For an undirected weighted graph with adjacency matrix \(A\), define the diagonal degree matrix
\[
D=\operatorname{diag}(d_1,\ldots,d_n),
\qquad d_i=\sum_j A_{ij}.
\]
The combinatorial graph Laplacian is
\[
\boxed{L=D-A.}
\]
The symmetric normalized Laplacian is
\[
L_{\mathrm{sym}}=I-D^{-1/2}AD^{-1/2}
\]
when all degrees are positive. The random-walk transition matrix is
\[
P=D^{-1}A.
\]
If \(x\in\mathbb R^n\) is a graph signal, then
\[
(Lx)_i=d_i x_i-\sum_j A_{ij}x_j
=\sum_j A_{ij}(x_i-x_j).
\]
Thus \(Lx\) compares each node to its neighbors.

\paragraph{Core theorem: Laplacian quadratic form.}
For an undirected weighted graph,
\[
\boxed{x^\top Lx=\frac12\sum_{i,j=1}^n A_{ij}(x_i-x_j)^2.}
\]
Therefore \(L\) is positive semidefinite, and \(x^\top Lx\) is small when \(x\) varies slowly across high-weight edges.

\emph{Proof.} Expand:
\[
x^\top Lx=x^\top Dx-x^\top Ax
=\sum_i d_i x_i^2-\sum_{i,j}A_{ij}x_i x_j.
\]
Since \(d_i=\sum_j A_{ij}\),
\[
\sum_i d_i x_i^2=\sum_{i,j}A_{ij}x_i^2.
\]
Because \(A_{ij}=A_{ji}\),
\[
\frac12\sum_{i,j}A_{ij}(x_i-x_j)^2
=\frac12\sum_{i,j}A_{ij}(x_i^2-2x_ix_j+x_j^2)
=\sum_i d_ix_i^2-\sum_{i,j}A_{ij}x_ix_j.
\]
This equals \(x^\top Lx\).

\paragraph{Concrete example.}
For the path \(1-2-3\),
\[
L=
\begin{pmatrix}
1&-1&0\\
-1&2&-1\\
0&-1&1
\end{pmatrix}.
\]
For \(x=(1,2,4)^\top\),
\[
x^\top Lx=(1-2)^2+(2-4)^2=1+4=5.
\]
The edge \(2-3\) contributes more because the signal jump is larger.

\paragraph{Computational neuroscience example.}
Suppose \(x_i\) is task-evoked activity in brain region \(i\), and \(A_{ij}\) measures structural connectivity. The quantity \(x^\top Lx\) is a graph-smoothness penalty: it is small when strongly connected regions have similar activity. This can be used as a regularizer, but it is a modeling assumption. Real activity can differ sharply across connected regions.

\paragraph{AI or machine-learning example.}
Semi-supervised learning on graphs often minimizes a loss plus a smoothness term
\[
\sum_{i\in\mathcal L}\ell(f_i,y_i)+\lambda f^\top Lf.
\]
The labeled nodes \(\mathcal L\) anchor the prediction, and the Laplacian term encourages neighboring nodes to have similar predictions. This works when graph edges connect examples likely to share labels; it fails when edges connect dissimilar examples.

\paragraph{Common misconceptions and failure modes.}
The Laplacian is not the adjacency matrix. The adjacency matrix says who is connected; the Laplacian compares each node to its neighbors. A small Laplacian energy means smoothness on the chosen graph, not correctness. Isolated nodes require care because normalized formulas use degree inverses. Directed graphs need different Laplacian conventions.

\paragraph{Exercises.}
\begin{enumerate}[label=\textbf{29.2.\arabic*.}, leftmargin=*]
\item \textbf{Mechanical.} Compute \(D\) and \(L\) for the triangle graph.
\item \textbf{Proof or derivation.} Prove that \(L\mathbf 1=0\), where \(\mathbf 1=(1,\ldots,1)^\top\).
\item \textbf{Numerical/computational.} For the four-node chain and \(x=(0,1,1,3)^\top\), compute \(x^\top Lx\) by summing squared edge differences.
\item \textbf{Interpretation.} What does a large value of \(x^\top Lx\) suggest about a graph signal on a connectome?
\item \textbf{Misconception/debugging.} A student says, ``Regularizing by \(x^\top Lx\) is always biologically correct for brain data.'' Give a reason this may be only a modeling assumption.
\end{enumerate}

\subsection{29.3 Spectral graph methods}

\paragraph{Guiding question.}
How do eigenvectors of the graph Laplacian reveal large-scale graph structure and graph-frequency modes?

\paragraph{Beginner-facing intuition.}
Fourier modes on a time line are eigenvectors of a shift-invariant operator. On a graph, the Laplacian eigenvectors play a similar role. Low-eigenvalue vectors vary slowly across edges. High-eigenvalue vectors oscillate rapidly across edges. This gives a graph version of frequency.

The second-smallest Laplacian eigenvector is especially important. It often separates a graph into two weakly connected parts and is called the Fiedler vector.

\paragraph{Formal definitions and notation.}
For an undirected graph, \(L\) is symmetric and positive semidefinite. It has an orthonormal eigenbasis
\[
Lu_k=\lambda_k u_k,\qquad
0=\lambda_1\leq\lambda_2\leq\cdots\leq\lambda_n.
\]
The graph Fourier coefficients of a signal \(x\in\mathbb R^n\) are
\[
\widehat x_k=u_k^\top x,
\qquad
x=\sum_{k=1}^n \widehat x_k u_k.
\]
Low values of \(\lambda_k\) correspond to smooth graph modes because
\[
u_k^\top L u_k=\lambda_k.
\]
If the graph has \(c\) connected components, then the nullspace of \(L\) has dimension \(c\). Its basis can be chosen as indicator vectors of the components.

\paragraph{Core mechanism: spectral relaxation of a cut.}
Partitioning a graph into two balanced groups is combinatorial. Spectral clustering relaxes the discrete assignment to a real-valued vector. The Rayleigh quotient
\[
R(x)=\frac{x^\top Lx}{x^\top x}
\]
is minimized by constant vectors with value \(0\), which do not split the graph. If we require \(x\perp\mathbf 1\), the minimizer is the eigenvector \(u_2\) with eigenvalue \(\lambda_2\):
\[
\min_{x\neq0,\ x\perp\mathbf 1}\frac{x^\top Lx}{x^\top x}=\lambda_2.
\]
Thresholding the entries of \(u_2\) gives a two-way spectral partition. This is an approximation to a discrete cut objective, not an exact solution in every graph.

\paragraph{Concrete example.}
For the four-node path \(1-2-3-4\), a Fiedler vector has entries approximately
\[
u_2\approx(-0.653,-0.271,0.271,0.653).
\]
The sign split separates \(\{1,2\}\) from \(\{3,4\}\), which matches the intuitive weakest middle cut of the chain. The vector varies slowly: neighboring entries are closer than distant entries.

\paragraph{Computational neuroscience example.}
In connectome analysis, Laplacian eigenvectors can reveal gradients across cortical regions. A low-frequency graph mode may vary smoothly over anatomical connectivity and summarize a large-scale axis of organization. This does not mean the brain literally computes that eigenvector; it means the eigenvector is a mathematical coordinate for structured variation on the network.

\paragraph{AI or machine-learning example.}
Spectral clustering embeds node \(i\) using entries of several low-frequency eigenvectors, then clusters those embedded points. Graph neural networks can also be read as graph-signal processors: repeated neighborhood averaging tends to emphasize low-frequency graph components, while more expressive architectures try to preserve task-relevant high-frequency information when labels alternate across edges.

\paragraph{Common misconceptions and failure modes.}
The smallest Laplacian eigenvector is usually constant on a connected graph, so it does not cluster. The Fiedler vector is useful but not infallible; degree heterogeneity, noise, and unbalanced communities can mislead it. Low graph frequency means smooth on the chosen graph, not slow in time. Spectral embeddings are coordinates chosen by a matrix, not direct physical locations unless the graph encodes physical geometry.

\paragraph{Exercises.}
\begin{enumerate}[label=\textbf{29.3.\arabic*.}, leftmargin=*]
\item \textbf{Mechanical.} For a connected graph, what is the eigenvalue associated with \(\mathbf 1\)?
\item \textbf{Proof or derivation.} Use the Rayleigh quotient to explain why the eigenvector with the second-smallest eigenvalue minimizes \(x^\top Lx\) among unit vectors orthogonal to \(\mathbf 1\).
\item \textbf{Numerical/computational.} For the two-component graph consisting of edges \(1-2\) and \(3-4\), find two independent vectors in the nullspace of \(L\).
\item \textbf{Interpretation.} What does it mean for a graph signal to have most of its energy in low Laplacian eigenmodes?
\item \textbf{Misconception/debugging.} A student says, ``The Fiedler vector always finds the true biological modules.'' Explain why the result depends on the graph construction and the chosen notion of module.
\end{enumerate}

\subsection{29.4 Network science for neuroscience and ML}

\paragraph{Guiding question.}
Which network summaries are useful in neuroscience and machine learning, and what assumptions do they carry?

\paragraph{Beginner-facing intuition.}
Network science studies patterns beyond individual edges: hubs, communities, clustering, paths, centrality, motifs, and small-world structure. These summaries can be useful, but each one compresses the graph. A centrality score, for example, is not a universal measure of importance; it measures importance according to a specific mathematical criterion.

\paragraph{Formal definitions and notation.}
The degree distribution records how many nodes have each degree. The local clustering coefficient of node \(i\) in an unweighted undirected graph is
\[
C_i=\frac{\text{number of edges among neighbors of }i}
{\binom{d_i}{2}},
\]
when \(d_i\geq2\). A shortest-path distance \(d_G(i,j)\) is the fewest edges in any path from \(i\) to \(j\). A community is an informal term for a set of nodes with many internal edges and fewer external edges; modularity is one common objective that formalizes this idea for a proposed partition.

A message-passing neural network updates node features by
\[
h_i^{(\ell+1)}
=\phi_\ell\!\left(h_i^{(\ell)},\ \operatorname{AGG}_{j\in N(i)}\psi_\ell(h_i^{(\ell)},h_j^{(\ell)},A_{ij})\right),
\]
where \(h_i^{(\ell)}\) is the feature vector of node \(i\) at layer \(\ell\), \(\operatorname{AGG}\) is a permutation-invariant aggregation such as sum or mean, and \(\phi_\ell,\psi_\ell\) are learned functions.

\paragraph{Core mechanism: random-walk diffusion on a graph.}
If \(d_i>0\), the random-walk matrix
\[
P=D^{-1}A
\]
has rows summing to one. A distribution row vector evolves by
\[
\pi_{t+1}=\pi_t P.
\]
A graph signal can also be smoothed by neighbor averaging:
\[
x_{t+1}=Px_t
\]
under a column convention adjusted consistently. Repeated averaging diffuses information over the network and tends to reduce sharp edge-to-edge differences. In continuous time, graph heat flow is
\[
\frac{dx}{dt}=-Lx,
\qquad
x(t)=e^{-tL}x(0).
\]
Laplacian eigenmodes decay as
\[
e^{-t\lambda_k}.
\]
High-\(\lambda_k\) modes decay faster, so heat flow is graph low-pass filtering.

\paragraph{Concrete example.}
Consider a node \(i\) with three neighbors. If those neighbors have two edges among themselves, then
\[
C_i=\frac{2}{\binom32}=\frac23.
\]
For a star graph with one center connected to \(m\) leaves, the center has degree \(m\) but clustering coefficient \(0\), because its neighbors are not connected to each other. Degree and clustering measure different aspects of network organization.

\paragraph{Computational neuroscience example.}
Structural connectomes use anatomical edges, functional networks often use statistical dependence between activity traces, and effective connectivity tries to model directed influence. These are different graph constructions. A hub in a structural connectome may not be a hub in a task-specific functional graph. Network summaries must be interpreted relative to the measurement and edge definition.

\paragraph{AI or machine-learning example.}
Graph neural networks use message passing to combine node features with neighborhood information. On a molecular graph, atoms pass messages along bonds. On a citation graph, papers aggregate information from cited or citing papers. Oversmoothing can occur when many layers make node representations too similar, echoing the graph diffusion mechanism above.

\paragraph{Common misconceptions and failure modes.}
Correlation networks do not prove direct causal interaction. Thresholding a dense correlation matrix can change graph metrics sharply. Community-detection algorithms can find modules in random graphs if the objective is forced. In GNNs, more message-passing layers do not always mean better use of graph structure; they can erase useful node differences.

\paragraph{Exercises.}
\begin{enumerate}[label=\textbf{29.4.\arabic*.}, leftmargin=*]
\item \textbf{Mechanical.} Compute the clustering coefficient of the center of a four-leaf star graph.
\item \textbf{Proof or derivation.} Show that \(P=D^{-1}A\) has rows summing to one when every degree is positive.
\item \textbf{Numerical/computational.} For the path \(1-2-3\), compute the one-step random-walk distribution starting from node \(2\).
\item \textbf{Interpretation.} Explain the difference between structural, functional, and effective connectivity in graph terms.
\item \textbf{Misconception/debugging.} A student says, ``A high correlation edge means one brain region drives the other.'' Explain why correlation alone does not establish directed causality.
\end{enumerate}

\paragraph{Closing bridge.}
Graphs add relational structure to the finite-dimensional mathematics of vectors and matrices. Adjacency matrices encode who is connected, Laplacians encode edge-wise differences, and eigenvectors provide graph-frequency coordinates. The next chapter adds action: controlled systems and dynamic programming ask how an agent or controller should move through a state space, often represented implicitly as a graph of possible transitions.

\clearpage
